%%=============================================================================
%% Voorwoord
%%=============================================================================

\chapter*{\IfLanguageName{dutch}{Woord vooraf}{Preface}}%
\label{ch:voorwoord}

%% TODO:
%% Het voorwoord is het enige deel van de bachelorproef waar je vanuit je
%% eigen standpunt (``ik-vorm'') mag schrijven. Je kan hier bv. motiveren
%% waarom jij het onderwerp wil bespreken.
%% Vergeet ook niet te bedanken wie je geholpen/gesteund/... heeft

Deze bachelorproef onderzoekt hoe misleidende technieken op streamingdiensten kunnen worden opgespoord en op een gebruiksvriendelijke manier aan jongvolwassenen kunnen worden gesignaleerd. Het onderwerp vormde voor mij een interessante combinatie van gebruikerservaring, interfaceontwerp en softwareontwikkeling.

\vspace{1em}

Voor ik aan de opleiding Toegepaste Informatica begon, volgde ik Grafische en Digitale Media, waar ik afstudeerde als grafisch ontwerper. Mijn interesse in UI/UX vormde dan ook een belangrijke motivatie om voor dit onderwerp te kiezen. Kort voor de start van mijn onderzoek kwam ik zelf in aanraking met interfaces die gebruikers proberen te beïnvloeden. Ook binnen mijn omgeving deden zich gelijkaardige ervaringen voor, onder andere bij het stopzetten van een Amazon Prime-abonnement. Hierdoor begon ik mij af te vragen hoe vaak dergelijke praktijken voorkomen. Dit vormde de aanleiding om mij verder te verdiepen in misleidende technieken, waarvoor later de term dark patterns bleek te worden gebruikt.

\vspace{1em}

Tijdens het onderzoek ontdekte ik dat dark patterns vaak subtieler zijn dan ik aanvankelijk dacht. Naarmate ik er meer over leerde, begon ik ze ook steeds vaker te herkennen in interfaces die ik dagelijks gebruik. Daarnaast bracht het onderzoek mij in aanraking met Android-appontwikkeling en AI, domeinen waarin ik vooraf weinig ervaring had. Het was bijzonder waardevol om nieuwe technische kennis op te doen en tegelijk mijn achtergrond in grafisch ontwerp te kunnen combineren met mijn opleiding Toegepaste Informatica.

\vspace{1em}

Tot slot wil ik mijn promotor, mevrouw L. De Mol, bedanken voor haar begeleiding, waardevolle feedback en steun tijdens het onderzoek. Ook mijn co-promotors, de heer G. De Bruyn en de heer S. Hostyn, wil ik bedanken voor hun advies en bereidheid om mee te werken aan interviews. Daarnaast bedank ik J. De Smedt, M. Adams, M. Dhondt, M. Maratta en S. Landtsheer voor hun deelname aan de interviews, en M. Adams, M. Dhondt en E. De Man voor hun deelname aan de gebruikerstesten. Met deze bachelorproef sluit ik mijn opleiding Toegepaste Informatica af en kijk ik trots terug op een bijzonder leerrijke periode.